% !TEX root = ../main.tex
\section{Re-measuring $F$ and $I$ in the New Domain}
\label{sec:axes}

\Cref{sec:sample} moved a \emph{diagnosis} from one driving side to the other. The more useful
question for deployment is whether the \emph{method} moves: can a handful of counterfactual cases in
the new domain characterise how the policies behave there? We repeat both interventions on the
\NumSGscenes{} NAVSIM Singapore near-pedestrian scenes and read, per scene, the two behavioural
displacements that correspond to $\hat v_F$ and $\hat v_I$ at the output: $F$, the plan displacement
caused by removing the target pedestrian, and $I$, that caused by the night-style perturbation.

\begin{figure}[t]
\centering
\includegraphics[width=\columnwidth]{figures/cases.pdf}
\caption{\textbf{Four right-hand-drive cases, and displacement against the separation it buys.}
(a--d)~Ego at the origin, the vulnerable road user's (VRU) logged future dashed from the star, and DiffusionDriveV2's own BEV
semantic head as the background; every policy receives the frame's logged navigation, so the turns are
commanded. (a)~Clearances
$0.3$--$0.9$\,m. (b)~SimLingo passes at $3.2$\,m, the other five at $6.7$--$10$\,m. (c,~d)~The
departing policies and their planned distances are named in the panel; Boston departure rates are
DiffusionDriveV2 81\%, AutoVLA 9\%, DiffusionDrive 2.8\%. (e)~Every policy--scene cell of the
\NumSGscenes{} Singapore scenes: $F$ against $\Delta S$; band: the 10--90\% spread of $\Delta S$ per
$F$ bin; dashed line: $F=0.5$\,m.}
\label{fig:cases}
\end{figure}

\paragraph{When the plan does move, which way does it move?}
$F$ measures how far a plan moves, not where it moves to. We therefore decompose it (\cref{fig:cases}e), reading $F$ against $I$ of
\eqref{eq:FI}---the same policy's jitter floor under an unrelated edit---and against $\Delta S$, the
change in clearance to the pedestrian's logged future. A response counts as avoidance only if all three hold: the plan moved
($F\ge0.5$\,m), it moved more than the policy's own jitter ($F>I$), and the movement bought distance
($\Delta S\ge0.5$\,m).

Three readings follow. First, the pedestrian outweighs the irrelevant edit in only
\NumFgtIlo--\NumFgtIhi\% of scenes, so most of what $F$ registers sits inside the policy's own
variability---the pinched band of \cref{fig:cases}e. Second, avoidance is rare where it is called for. The logged ego is
stationary through most of this corpus, so a plan that does not move is usually the correct one; we
therefore replay every scene at input speeds of 4 and 8\,m/s, as in \cref{sec:checkup}, and read only the
cells in which the pedestrian-blind plan does come within 1\,m of the pedestrian's logged future.
Across \NumCfModels{} policies that is \NumCfNeed{} of \NumCfCells{} cells (\NumCfNeedLo--\NumCfNeedHi\%
per policy), and \NumCfReal{} of them---\NumCfRealPct\%---are genuine avoidance. On those cells the
median clearance changes by \NumCfDsLo--\NumCfDsHi\,m and the median planned distance by
\NumCfArcLo--\NumCfArcHi\,m: facing a pedestrian it would otherwise reach, no policy diverts and none
slows. Third, the displacement that does occur is not aimed: of the
\NumMoveN{} cells in which the plan moved, \NumMoveFar{} increased the clearance and \NumMoveNear{}
decreased it. Per policy the split differs sharply: DiffusionDrive reacts in few
scenes but in the right direction (median $\Delta S=\NumDDmed$\,m), while SimLingo moves in more
scenes than the other five combined and moves \emph{towards} the pedestrian more often than away
(\NumSLnear{} against \NumSLfar, median $\NumSLmed$\,m); it also has the highest rate of sub-1.5\,s
time-to-proximity (\cref{tab:bench}). Finally, what a policy yields does not grow with the hazard: the
rank correlation between $\Delta S$ and the deceleration the scene demands is
\NumRhoLo--\NumRhoHi{}.

An exam that rewarded responsiveness alone would therefore rank the least safe policy first.

\paragraph{What a small sample in the new domain buys}
We draw $k$ scenes at random from the pool, re-measure $F$ and $I$ on them and compare against the
large-sample values, 200 draws per $k$. Ordering and level behave very differently. The ordering of
how dangerously the policies behave---the share of scenes whose minimum time-to-proximity falls below
$1.5$\,s---is recovered almost immediately: $\rho={\NumBehRhoTen}$ at $k=10$ and $\rho={\NumBehRhoTwenty}$
at $k=20$ ($\rho\ge0.6$ in \NumBehHitTwenty\% of draws), rising to $\rho={\NumBehRhoForty}$ at $k=40$.
The level is not: the same $k=20$ sample misstates the violation rate by \NumBehErrTwenty{} points and
$\CFR$ by $\NumCfrErrTwenty$; the rule hit-rates sit in between, at \NumPFErrTwenty{} points. Importing
the Boston $\CFR$ values instead misstates the Singapore ones by $\NumLhdCfrErr$, already worse than
re-measuring $k=40$ fresh scenes ($\NumCfrErrForty$), while 80 scenes reach $\NumCfrErrBig$ and 160
$\NumCfrErrHuge$. Values measured elsewhere are worth about forty scenes measured here.

\paragraph{Which exam travels}
Of the exams in the report, the one that predicts the new domain's behaviour is not the one that
measures how \emph{much} a policy reacts, but the one that measures how \emph{steadily} it drives.
Correlating each Boston exam against the Singapore rate of sub-1.5\,s time-to-proximity over the six
policies, specificity gives $\rho={\NumSpBeh}$---the only value significant at the two-sided 5\% level
for $n{=}6$---against ${\NumExpBeh}$ for exposure and ${\NumCfrBeh}$ for $\CFR$. Specificity asks how far a plan moves when an
\emph{irrelevant} part of the scene is edited; a policy that answers ``not at all'' drives steadily,
and steadiness is what survives the change of side. (Exposure travels too: on the larger
original-image set of \cref{sec:bench} it reaches $0.94$ against the same target. Specificity is the
stronger predictor and the only one that clears significance here.) The decomposition above says why:
displacement that is not aimed away from the pedestrian is not protection, and a policy producing a
lot of it is dangerous on both sides.
